\documentclass[10pt]{article}
\providecommand{\FIGDIR}{fig}
\newcommand{\DIGESTDATE}{15 August 2026}
% ===================== PM Research Watch - shared preamble =====================
\usepackage[T1]{fontenc}
\usepackage[utf8]{inputenc}
\usepackage{amsmath,amssymb}
\usepackage{tgpagella}
\usepackage{tgheros}
\renewcommand{\sfdefault}{qhv}
\renewcommand{\ttdefault}{lmtt}
\usepackage{microtype}
\usepackage[english]{babel}

\usepackage[a4paper,top=17mm,bottom=17mm,left=17mm,right=17mm,
            headsep=6mm,footskip=10mm]{geometry}
\usepackage{graphicx}
\usepackage{xcolor}
\usepackage{colortbl}
\usepackage{tikz}
\usetikzlibrary{calc}
\usepackage[most]{tcolorbox}
\usepackage{titlesec}
\usepackage{fancyhdr}
\usepackage{booktabs}
\usepackage{array}
\usepackage{longtable}
\usepackage{enumitem}
\usepackage{multicol}
\usepackage{ragged2e}
\usepackage[hidelinks]{hyperref}

% ------------------------------------------------------------------ palette
\definecolor{Ink}    {HTML}{14262E}
\definecolor{Deep}   {HTML}{0F4C5C}
\definecolor{Teal}   {HTML}{1F7A8C}
\definecolor{Sky}    {HTML}{4FA3C4}
\definecolor{Sage}   {HTML}{6FA287}
\definecolor{Amber}  {HTML}{E3A008}
\definecolor{Clay}   {HTML}{B5643C}
\definecolor{Coral}  {HTML}{D1495B}
\definecolor{Violet} {HTML}{6A5B9E}
\definecolor{Slate}  {HTML}{4A6572}
\definecolor{Paper}  {HTML}{FBF9F5}
\definecolor{Mist}   {HTML}{EEEAE2}
\definecolor{Rule}   {HTML}{D9D3C9}

\pagecolor{Paper}
\color{Ink}

\hypersetup{colorlinks=true, linkcolor=Deep, urlcolor=Teal, citecolor=Teal,
            pdftitle={PM Research Watch}, pdfauthor={Automated daily digest}}

% ------------------------------------------------------------------ headings
\titleformat{\section}
  {\sffamily\bfseries\large\color{Deep}}{}{0pt}
  {}[{\vspace{1.2mm}\color{Amber}\titlerule[1.1pt]}]
\titlespacing*{\section}{0pt}{5.5mm}{2.8mm}

\titleformat{\subsection}
  {\sffamily\bfseries\normalsize\color{Teal}}{}{0pt}{}
\titlespacing*{\subsection}{0pt}{3.6mm}{1.4mm}

\setlist[itemize]{leftmargin=4.2mm, itemsep=1.1mm, topsep=1.3mm, parsep=0pt,
                  label=\textcolor{Teal}{\footnotesize$\blacktriangleright$}}

% ------------------------------------------------------------------ header/footer
\pagestyle{fancy}
\fancyhf{}
\renewcommand{\headrulewidth}{0pt}
\renewcommand{\footrulewidth}{0pt}
\fancyhead[L]{\sffamily\scriptsize\color{Slate}PM RESEARCH WATCH\;\textcolor{Rule}{\textbar}\;\DIGESTDATE}
\fancyhead[R]{\sffamily\scriptsize\color{Slate}Particulate matter monitoring \& health effects}
\fancyfoot[C]{\sffamily\scriptsize\color{Slate}\thepage}

% ------------------------------------------------------------------ components
% metric tile
\newcommand{\metric}[3]{%
  \begin{tcolorbox}[enhanced, colback=white, colframe=Rule, boxrule=0.5pt,
    arc=1.6mm, left=2.2mm, right=2.2mm, top=1.8mm, bottom=1.8mm,
    borderline west={1.7mm}{0pt}{#3}, nobeforeafter, valign=center,
    height=17mm, width=\linewidth]
    \RaggedRight
    {\sffamily\bfseries\fontsize{17}{17}\selectfont\color{#3}#1}\\[0.5mm]
    {\sffamily\fontsize{6.5}{7.4}\selectfont\color{Slate}#2}
  \end{tcolorbox}}

% pull-quote / key insight box
\newtcolorbox{keybox}[1][]{enhanced, colback=Mist, colframe=Mist,
  arc=1.6mm, left=3.2mm, right=3.2mm, top=2.4mm, bottom=2.4mm,
  borderline west={1.7mm}{0pt}{Amber}, #1}

% subtopic band
\newcommand{\band}[2]{%
  \vspace{3.2mm}\noindent
  \begin{tcolorbox}[enhanced, colback=#2, colframe=#2, arc=1.2mm,
    left=2.6mm, right=2.6mm, top=1.1mm, bottom=1.1mm, boxrule=0pt,
    width=\linewidth, nobeforeafter]
    {\sffamily\bfseries\small\color{white}#1}
  \end{tcolorbox}\vspace{1.6mm}}

% one paper entry: {short cite}{title}{journal}{doi}{body}
\newcommand{\paperentry}[5]{%
  \par\vspace{1.9mm}\noindent
  {\sffamily\bfseries\footnotesize\color{Ink}#1}\;%
  {\sffamily\fontsize{7}{8}\selectfont\color{Slate}#3}\\[0.4mm]
  {\footnotesize\itshape\color{Deep}#2}\\[0.6mm]
  {\footnotesize #4}\\[0.3mm]
  {\fontsize{6.6}{7.6}\selectfont\color{Slate}\texttt{doi:}\,\href{https://doi.org/#5}{\texttt{\detokenize{#5}}}}
  \par}

% caption style
\newcommand{\figcap}[1]{\par\vspace{1.0mm}
  {\fontsize{7.2}{8.6}\selectfont\color{Slate}\textsf{\textbf{Fig.}}\; #1\par}\vspace{1.5mm}}

\providecommand{\permil}{\textperthousand}

\setlength{\parindent}{0pt}
\setlength{\parskip}{0pt}
\renewcommand{\arraystretch}{1.18}


\begin{document}

% ==================================================================== MASTHEAD
\thispagestyle{empty}
\begin{tikzpicture}[remember picture, overlay]
  \fill[Deep] (current page.north west) rectangle
      ($(current page.north east)+(0,-67mm)$);
  \fill[Amber] ($(current page.north west)+(0,-67mm)$) rectangle
      ($(current page.north east)+(0,-68.4mm)$);
  \foreach \x/\y/\r in {22/12/0.9, 41/26/1.5, 63/9/0.7, 88/31/1.1, 112/17/1.9,
                        138/29/0.8, 159/13/1.3, 178/54/1.0, 196/21/1.6,
                        30/44/0.6, 74/57/1.2, 125/48/0.9, 168/6/0.8, 191/59/0.7,
                        52/18/1.0, 100/60/0.6, 148/20/1.1, 66/49/1.4, 205/45/0.8}
    \fill[Sky, opacity=0.30] ($(current page.north west)+(\x mm,-\y mm)$) circle (\r mm);
\end{tikzpicture}

\vspace*{4mm}
{\sffamily\fontsize{7.6}{9}\selectfont\color{Sky}%
 AUTOMATED DAILY LITERATURE DIGEST\;\textbullet\;PhD RESEARCH WATCH\par}
\vspace{2.2mm}
{\sffamily\bfseries\fontsize{27}{29}\selectfont\color{white}%
 Particulate Matter\par}
{\sffamily\fontsize{27}{29}\selectfont\color{Sky}%
 Monitoring \& Health Effects\par}
\vspace{5.0mm}
{\sffamily\fontsize{8.4}{10}\selectfont\color{white}%
 Issue 2026-08-15\;\;\textcolor{Amber}{\textbar}\;\;Entry window 15 August 2026\;\;%
 \textcolor{Amber}{\textbar}\;\;21 records carried, 26 rejected\par}
\vspace{18mm}

% -------------------------------------------------------------- metric strip
\noindent
\begin{minipage}[t]{0.190\linewidth}\metric{21}{RECORDS CARRIED\\(26 REJECTED)}{Deep}\end{minipage}\hfill
\begin{minipage}[t]{0.190\linewidth}\metric{9}{SUBTOPIC CLUSTERS\\REPRESENTED}{Teal}\end{minipage}\hfill
\begin{minipage}[t]{0.190\linewidth}\metric{5}{EFFECT ESTIMATES\\WITH INTERVALS}{Sky}\end{minipage}\hfill
\begin{minipage}[t]{0.190\linewidth}\metric{7}{INSTRUMENTATION /\\EXPOSURE RECORDS}{Amber}\end{minipage}\hfill
\begin{minipage}[t]{0.190\linewidth}\metric{2}{METADATA-ONLY\\(NO ABSTRACT)}{Clay}\end{minipage}

\vspace{6mm}

% -------------------------------------------------------------- signal of day
\section*{\sffamily\bfseries\color{Deep}Signal of the day}
\vspace{-1.5mm}{\color{Amber}\rule{\linewidth}{1.1pt}}\vspace{2.5mm}

\begin{keybox}
{\sffamily\bfseries\small\color{Deep}Two large cohorts extend PM$_{2.5}$ dose--response
into South Asia and sub-Saharan Africa on the same day, and both rest entirely on modelled
exposure in the regions where the models are least validated.}\\[1.4mm]
\footnotesize
Mandal et al.\ (\textit{Kidney Int Rep}) follow 12{,}271 adults in Chennai and Delhi across
three eGFR measurements and report \textbf{$-$0.32 (95\% CI $-$0.49 to $-$0.15) and
$-$0.42 ($-$0.57 to $-$0.27)\,ml/min/1.73\,m$^2$ per 5\,$\mu$g/m$^3$} of annual PM$_{2.5}$.
Moro et al.\ (\textit{Environ Res}) find \textbf{OR 1.06 (1.02--1.11) per
1\,$\mu$g/m$^3$ for hypertension} in 1{,}833 rural Ghanaian adults, PAF 18.6\%, with no
association for metabolic syndrome or any other component.

Both are worth having --- renal and sub-Saharan cardiometabolic endpoints are genuine
evidence gaps. But the exposure in Mandal is a 1\,km national spatiotemporal model at
Delhi annual means of \textbf{118--130\,$\mu$g/m$^3$}, and in Moro a satellite-derived
model over rural Ghana. Neither region has the reference-monitor density that calibrated
the modelling approaches these estimates inherit, and neither paper reports out-of-sample
model error at the exposure levels it is used at. A per-unit odds ratio of 1.06 is
recoverable only if the exposure contrast is measured to better than a unit. That is a
measurement claim, not an epidemiological one, and it is unaddressed in both papers ---
which is precisely the argument for dense low-cost networks in exactly these settings.
\end{keybox}

\vspace{3mm}

% -------------------------------------------------------------- exec brief
\section*{\sffamily\bfseries\color{Deep}What changed today}
\vspace{-1.5mm}{\color{Amber}\rule{\linewidth}{1.1pt}}\vspace{2.5mm}

\footnotesize
\begin{itemize}
\item \textbf{A community-run mobile LCS campaign reports its own year-to-year model
instability.} Cho et al.\ ran two summers of community-scientist mobile PM$_{2.5}$
monitoring in Norwood, Ohio, then fitted street-level land-use regression. Cross-validated
$R^2$ went from \textbf{0.17--0.25 in 2020 to 0.55--0.61 in 2021} on the same streets with
the same protocol. The paper reads the difference as pandemic-phase source recovery; it is
equally consistent with 2020 having too little exposure contrast for LUR to resolve.
Either way, publishing both years rather than the good one is the right call.

\item \textbf{Combined cold-spell and heavy-pollution exposure is supra-additive for
circulatory mortality.} Zhang et al.\ pooled 70 districts across the Beijing--Tianjin--Hebei
region, cold seasons 2013--2018. \textbf{CRR 1.35 (1.14--1.60)} for the overlapping event
against \textbf{1.14 (1.09--1.20)} for cold spell alone and \textbf{1.08 (1.00--1.16)} for
heavy pollution alone --- the last of which only touches unity. Two lagged
single-exposure warnings would each have under-called the combined day.

\item \textbf{A continuous HNO$_3$/particulate-nitrate method at sub-0.15\,$\mu$g/m$^3$
detection limits.} Alexa et al.\ pair a wet diffusion denuder with a continuous aerosol
sampler and chemiluminescence detection, recording at 1\,s intervals, with LODs of
\textbf{0.13\,$\mu$g/m$^3$ for HNO$_3$ and 0.10\,$\mu$g/m$^3$ for NO$_3^-$}, validated
against filter/denuder reference sampling across Brno winter and summer campaigns. Nitrate
partitioning is exactly the fast process filter methods average away and low-cost optical
sensors cannot see at all.

\item \textbf{Twelve French sites, seven years, one harmonised PMF --- and the CTM is
systematically biased.} Chebaicheb et al.\ apply harmonised rolling PMF to ACSM organic
aerosol at 12 (sub)urban sites. Primary OA factors are $\approx$33\% of OA mass; oxygenated
OA dominates everywhere; wintertime LO-OOA tracks BBOA, implicating residential biomass
burning in secondary formation. Comparison against CHIMERE shows systematic component-level
bias. This is the reference dataset any speciated-PM epidemiology in France should be built
against.

\item \textbf{Nine of twenty-one records carry no health endpoint.} At 43\% that share is
flat against 12--14 August (46\%, 41\%, 43\%) and the split is now a stable property of the
corpus rather than a daily fluctuation. Every \textit{Sensing} and \textit{Exposure} record
measures particles without measuring a person; every cohort record models exposure without
measuring particles. Unlike 14 August, no record today measured both on the same subjects.
\end{itemize}

\clearpage

% ==================================================================== ANALYTICS
\section*{\sffamily\bfseries\color{Deep}Corpus analytics}
\vspace{-1.5mm}{\color{Amber}\rule{\linewidth}{1.1pt}}\vspace{3mm}

\noindent\includegraphics[width=\linewidth]{\FIGDIR/f1_subtopics.png}
\figcap{Subtopic assignment of the 21 carried records, single-label by dominant
contribution. Nine of ten clusters are occupied and \textit{Neuro / mental health} is empty
for the second consecutive issue --- and on 11 of the 21 issues to date, including an
unbroken six-day run from 3 to 9 August, so this is the watch's least reliable cluster
rather than a new gap.
\textit{Sensing, forecasting \& instrumentation} leads at 4; \textit{Mechanistic
toxicology}, \textit{Cardiovascular \& metabolic}, \textit{Exposure assessment \&
modelling} and \textit{Burden, policy \& mitigation} each hold 3.}

\vspace{2mm}
\noindent\includegraphics[width=\linewidth]{\FIGDIR/f2_design_endpoint.png}
\figcap{Left --- study architecture. \textit{Measurement campaign} (6) and \textit{Modelling
/ inventory} (5) together take eleven of twenty-one records, the most measurement-weighted
day of the week; observational designs split 3 cohort, 2 acute, 1 cross-sectional, and
there is no \textit{Trial / intervention} record at all. Right --- health endpoint.
\textbf{No health endpoint} is the largest bar at 9, covering the four sensing records, the
three exposure/modelling records, the indoor diesel-tracer campaign and the ship-emission
inventory. Cardiovascular endpoints follow at 4.}

\vspace{2mm}
\noindent\includegraphics[width=\linewidth]{\FIGDIR/f3_metric_geography.png}
\figcap{Left --- particle metric. PM$_{2.5}$-only accounts for 11 records. Three records
depart from mass: Ranjbar et al.'s nitro-PAH speciation, Chebaicheb et al.'s ACSM organic
factors (binned with unspeciated/emissions), and Kafi et al., whose Sentinel-5P evaluation
carries \emph{no surface PM product at all} and is binned by the classifier as a composite
index --- read that slice as an absence, not a measurement. Right --- geography. China
contributes 6 records, Europe 4, North America and East Asia (ex-China) 3 each; South Asia,
sub-Saharan Africa and MENA one apiece, and all three of those are the day's cohort or
satellite-evaluation records rather than instrumentation.}

\vspace{2mm}
\noindent\includegraphics[width=\linewidth]{\FIGDIR/f6_lifecourse.png}
\figcap{Life-course windows addressed. Counts are non-exclusive and describe this issue
only. Preconception/in-utero holds three records --- its highest count in eight issues
(late-pregnancy anthropometry, gestational TiO$_2$ inhalation, and a semen-quality cohort
whose exposure window is preconceptional). Puberty and adolescence is empty, as on three of
the previous five issues.}

\vspace{2mm}
\noindent\includegraphics[width=\linewidth]{\FIGDIR/f4_heatmap.png}
\figcap{Subtopic $\times$ study-architecture cross-tabulation. Twenty-one records occupy
13 of the 70 available cells and no cell holds more than 3. The silo is visible as a clean
block structure rather than a diagonal: \textit{Sensing} and \textit{Exposure} sit entirely
in the \textit{Measurement campaign} and \textit{Modelling / inventory} columns,
\textit{Mechanistic toxicology} entirely in \textit{Experimental / toxicology}, and every
human-endpoint subtopic in an observational column --- with one exception, the
\textit{Cardiovascular} cell in \textit{Review / synthesis}, which is the metadata-only
methodological letter. Unlike 14 August, there is no record this issue in which exposure and
outcome were measured by the same group on the same subjects.}

\vspace{2mm}
\noindent\includegraphics[width=\linewidth]{\FIGDIR/f5_forest.png}
\figcap{The 5 interval estimates carried today, on a common logarithmic axis --- the
thinnest forest plot of the week, because most of today's corpus is instrumentation.
\textbf{The three Zhang et al.\ estimates are not independent of one another}: they are
nested contrasts from the same 70-district meta-analysis, and the combined-exposure
estimate shares person-time with both single-exposure estimates. Shi et al.\ report a
99th-percentile-versus-median contrast at cumulative lag 0--5\,d, and Moro et al.\ a
per-1\,$\mu$g/m$^3$ odds ratio, so the exposure contrasts differ by more than the plot can
show. Mandal et al.'s eGFR slopes are on an additive ml/min/1.73\,m$^2$ scale and cannot be
placed on a log ratio axis; they are reported in the digest and the signal box instead.}

\clearpage

% ==================================================================== DIGEST
\section*{\sffamily\bfseries\color{Deep}Paper-level digest}
\vspace{-1.5mm}{\color{Amber}\rule{\linewidth}{1.1pt}}\vspace{1mm}

{\fontsize{7.4}{9}\selectfont\color{Slate}Ordered by cluster. Each entry gives the
methodological core and the finding that matters, not an abstract paraphrase.\par}

\band{Sensing, forecasting \& instrumentation\hfill 4 records}{Amber}

\paperentry{Cho et al.\ 2026}
{Greenspace, land use and PM$_{2.5}$ across pandemic phases: a community scientist-led mobile sensing campaign}
{Sci Total Environ}
{Community scientists ran low-cost mobile PM$_{2.5}$ monitoring across Norwood, Ohio, for
two summer periods --- the initial pandemic phase (2020) and the reopening phase (2021) ---
with a meteorology-adjusted difference-in-differences framework for the year contrast and
street-level linear land-use regression for spatial structure. \textbf{Higher 2021
concentrations were not fully explained by meteorology.} In the 2021 afternoon model,
PM$_{2.5}$ was elevated within 50\,m of office and manufacturing land use and inversely
associated with NDVI within 100\,m. \textbf{Cross-validated $R^2$ 0.17--0.25 (2020) rising
to 0.55--0.61 (2021).} The paper attributes the gap to source recovery, but a
low-contrast summer is an equally good explanation and is not tested. Health-survey scores
by ward are reported alongside but are not linked to measured exposure at the individual
level. The transferable part is the QC and sampling protocol for neighbourhoods with no
regulatory coverage, which is stated explicitly as the deliverable.}
{10.1016/j.scitotenv.2026.182124}

\paperentry{Alexa et al.\ 2026}
{Simultaneous continuous determination of gaseous HNO$_3$ and particulate nitrate by chemiluminescence}
{Talanta}
{Gaseous HNO$_3$ is sampled by wet diffusion denuder and particulate NO$_3^-$ in parallel
by continuous aerosol sampler, both analysed online by continuous flow analysers with
chemiluminescence detection --- no preconcentration step, signal recorded at 1\,s
intervals. Detection limits (3\,S/N) are \textbf{5.1\,nM and 12.8\,nM in solution,
corresponding to 0.13\,$\mu$g/m$^3$ (49.7\,ppt) for HNO$_3$ and 0.10\,$\mu$g/m$^3$ for
NO$_3^-$}. Applied to urban air in Brno across winter and summer campaigns 2022--2023 and
compared against filter sampling with NaF-coated dry denuders. The gas--particle
partitioning of nitrate is temperature- and humidity-driven on timescales of minutes, which
filter-based methods average away entirely and optical low-cost sensors cannot resolve at
all; a method at this time resolution is the reference any nitrate-aware correction scheme
would need. The abstract reports no field intercomparison statistics against the reference
method, only that the comparison was made.}
{10.1016/j.talanta.2026.130457}

\paperentry{Kafi et al.\ 2026}
{Evaluating the suitability of Sentinel-5P data for air quality monitoring in Tehran}
{Sci Total Environ}
{Six years (2019--2024) of TROPOMI Level-2 CO, NO$_2$, SO$_2$, O$_3$ and UV Aerosol Index
against 22 ground stations in a megacity with complex topography, assessed by Pearson and
Spearman correlation and Kling--Gupta Efficiency, plus a spatiotemporal trajectory
comparison. \textbf{Only NO$_2$ showed strong agreement; CO, SO$_2$ and O$_3$ were weak},
and the trajectory analysis diverged as well. The authors are explicit that UVAI represents
column aerosol loading rather than surface particulate concentration and warn against
reading it as PM --- a caveat routinely ignored in the satellite-PM literature and worth
citing for that reason alone. The diagnosis (coarse pixel size against a topographically
constrained urban basin) is plausible but not separated from the competing explanation of
station siting, and no vertical-profile or boundary-layer-height information is brought in
to test it.}
{10.1016/j.scitotenv.2026.182069}

\paperentry{Aparicio et al.\ 2026}
{Sensor system design and deployment for environmental and visitor monitoring in cultural heritage sites}
{Sensors (Basel)}
{Design, build and field validation of low-power low-cost nodes across five pilot sites in
four European countries under the EU ARGUS project, in two configurations: static nodes at
fixed positions and dynamic nodes carried by a quadruped robot or drone. Payload spans
thermal--humidity probes, volumetric water content, wind, pollution monitors and two
generations of visitor-counting devices, with Baltan\'as used as the iterative prototyping
site before wider deployment. Carried for the mobile-platform architecture --- a
sensor-on-robot deployment pattern that is directly transferable to indoor and
semi-enclosed PM mapping --- rather than for any air-quality result. \textbf{The abstract
reports no PM performance figures, no reference co-location and no calibration procedure},
which is the whole question for a low-cost node, so nothing here constrains data quality.
Tier C on that basis.}
{10.3390/s26154941}

\band{Exposure assessment \& modelling\hfill 3 records}{Amber}

\paperentry{Chebaicheb et al.\ 2026}
{A phenomenology of submicron organic aerosol components in urban France}
{Environ Int}
{The first nationwide harmonised organic-aerosol source apportionment for France:
time-resolved Aerosol Chemical Speciation Monitor measurements at 12 (sub)urban sites over
one to seven years each, with a single rolling PMF protocol applied under common
constraints and acceptance criteria across all datasets. Primary factors --- hydrocarbon-like,
biomass-burning and cooking-like OA --- account for \textbf{$\approx$33\% of OA mass};
oxygenated OA dominates at every site and splits into less- and more-oxidised components
everywhere except Poitiers and Strasbourg. Site-specific factors resolved a mixed
shipping/industrial component at Marseille Longchamp and an amine-related component at
Strasbourg and Creil. Wintertime LO-OOA correlates with BBOA, implicating residential
combustion in secondary as well as primary mass; summer LO-OOA tracks biogenic precursors.
\textbf{CHIMERE shows systematic component-level bias against the observations.} The
harmonisation is the contribution --- rolling PMF applied inconsistently is why the existing
European OA literature does not pool --- and the limitation is that (sub)urban sites only
cannot constrain the rural background these factors are transported through.}
{10.1016/j.envint.2026.110453}

\paperentry{Sui et al.\ 2026}
{Contrasting annual variation of organic and elemental carbon in the Pearl River Delta}
{J Environ Manage}
{Three years of continuous OC/EC observation (2019--2021) spanning the COVID-19 lockdown,
with PMF source apportionment. OC declined steadily, driven by suppressed secondary organic
carbon formation in the PRD's NOx-rich regime --- a relationship the authors report as
\emph{reversing} during lockdown when NOx fell sharply, which is the interesting result and
the one that identifies the regime. EC was flat or slightly rising, so fossil combustion
from industry and mobile sources has not responded to the control programme.
\textbf{Fossil-related emissions are $\approx$70\% of carbonaceous aerosol} and that share
did not fall over the period. A single-site three-year record cannot separate emission
change from meteorological variability, and no such adjustment is described; the lockdown
period is doing most of the causal work and is a one-off natural experiment with a
confounded meteorology.}
{10.1016/j.jenvman.2026.130708}

\paperentry{You et al.\ 2026}
{Satellite-based estimation of ambient ammonia and emission inventory uncertainty in South Korea}
{Environ Pollut}
{Monthly ground-level NH$_3$ at 15\,km resolution for 2013--2017, from Cross-track Infrared
Sounder retrievals refined by a two-stage framework: a linear mixed-effects model carrying
meteorology plus the Clean Air Policy Support System inventory, then a generalised additive
model over the LMM residuals. \textbf{Cross-validated $R^2$ = 0.73, MAE 0.15, RMSE 0.21}.
Concentrations peak in agricultural areas March--June; temperature is the dominant
meteorological control (\textbf{+18.9\% per +1\,$^{\circ}$C in June}), with relative
humidity and wind speed strongest in March. The GAM residual surface identifies
agricultural hotspots the inventory does not contain, which is the paper's actual claim ---
inventory gaps, not model skill. NH$_3$ matters here as the limiting precursor for ammonium
nitrate and sulfate, so a 15\,km monthly NH$_3$ field is a plausible predictor for
secondary PM$_{2.5}$ models, but the resolution is far coarser than the sources, and
validation against sparse ground NH$_3$ is circular with the sparsity that motivated the
work.}
{10.1016/j.envpol.2026.128903}

\band{Mechanistic toxicology\hfill 3 records}{Violet}

\paperentry{Mahmutovic Persson et al.\ 2026}
{Chronic aerosolized PM$_{2.5}$, plasma proteins and coronary S1P response in rats}
{Environ Res}
{Well-characterised urban PM$_{2.5}$ was re-aerosolised and delivered to rats under two
regimes designed to separate an acute high dose from a chronic lower dose at comparable
cumulative burden, with coronary vasomotor tone measured ex vivo by wire myograph and
systemic effects read from plasma sphingosine-1-phosphate and endothelin-1 plus
immunohistochemistry. \textbf{Only the chronic regime moved anything}: altered S1P vasomotor
response in coronary arteries and significantly decreased plasma total protein and ET-1.
The design is the point --- most PM inhalation toxicology uses acute high-dose exposure
because it is tractable, and this says that route can miss the effect entirely. Re-aerosolised
collected PM is not fresh ambient aerosol (semivolatiles are lost, morphology changes), the
ET-1 \emph{decrease} is not obviously in the direction the vascular-dysfunction hypothesis
predicts and is not reconciled, and no n or dose figures appear in the abstract.}
{10.1016/j.envres.2026.125391}

\paperentry{Seymore et al.\ 2026}
{Gestational nanoparticle inhalation alters placental glucose transporters and increases glucose flux}
{Toxicol Appl Pharmacol}
{Pregnant Sprague Dawley rats exposed by whole-body inhalation to TiO$_2$ nanoparticles at
\textbf{9.8 $\pm$ 1.0\,mg/m$^3$}, an occupational-level concentration, from gestational day
5 to 19; GLUT1/3/4 assayed in term placenta at GD20 and ex vivo placental perfusion used as
a functional test. Glut3 mRNA and GLUT1 protein fell, but GLUT4 expression and membrane
localisation of both GLUT1 and GLUT4 rose, and perfusion confirmed \textbf{increased}
glucose flux (AUC 95\% CI 77.4--127.5 exposed vs 39.1--73.6 control). The authors state
plainly that this contradicts their hypothesis. Read as a compensatory adaptation rather
than a benign finding --- upregulated placental glucose delivery is not obviously protective
in a fetal-growth-restriction model. The dose is roughly three orders of magnitude above
ambient PM$_{2.5}$ and the particle is a single engineered material, so this constrains
occupational nanomaterial exposure and says little directly about ambient PM.}
{10.1016/j.taap.2026.117996}

\paperentry{Chen et al.\ 2026}
{Gut microbiota signatures reflect PM$_{2.5}$-exacerbated asthma in a murine model}
{Allergy}
{\textbf{Metadata-only: a research letter deposited without an abstract, so nothing below the
title is verifiable from the record.} Registered here because the gut--lung axis in
PM-exacerbated allergic airway disease is a live mechanistic question and because the author
group (Academia Sinica Air Quality Research Center with NTU clinical immunology) pairs
characterised ambient aerosol with a murine model, which is unusual. No finding, effect size
or exposure characterisation is claimed by this issue, and the record supports no figure or
statistic in the analytics. Flagged for retrieval once the full text is indexed.}
{10.1111/all.70485}

\band{Cardiovascular \& metabolic\hfill 3 records}{Teal}

\paperentry{Zhang et al.\ 2026}
{Combined cold-spell and heavy-pollution exposure and circulatory mortality, Beijing--Tianjin--Hebei}
{Zhonghua Yu Fang Yi Xue Za Zhi}
{Daily mortality, PM$_{2.5}$ and meteorology from 70 districts and counties, cold seasons
(November--March) 2013--2018, mean daily PM$_{2.5}$ \textbf{98.58 $\pm$ 78.19\,$\mu$g/m$^3$}.
Three combined-exposure scenarios were constructed --- overlapping cold spell and heavy
pollution (A), a standalone pollution event within 7\,d of a cold spell ending (B), and a
standalone cold spell within 7\,d of a pollution event ending (C) --- with district-level
DLNMs pooled by random-effects meta-analysis. The overlapping event gave \textbf{CRR 1.35
(1.14--1.60)} against \textbf{1.14 (1.09--1.20)} for cold spell alone and \textbf{1.08
(1.00--1.16)} for heavy pollution alone. Supra-additivity is the substantive result, and it
is the case for joint rather than single-hazard alerting. Heavy pollution is defined by a
concentration threshold, not by composition, so the interaction could equally be a
composition effect: winter heating episodes in this region differ chemically from other
high-mass days, and the design cannot separate the two.}
{10.3760/cma.j.cn112150-20260105-00006}

\paperentry{Moro et al.\ 2026}
{Long-term PM$_{2.5}$, hypertension and metabolic syndrome in rural Ghana}
{Environ Res}
{Cross-sectional analysis of 1{,}833 AWI-Gen participants aged 40--60 in rural Ghana, with
2007--2015 mean PM$_{2.5}$ from a satellite-based model and MetS by harmonised criteria,
under sequentially adjusted multivariable logistic regression. \textbf{OR 1.06 (1.02--1.11)
per 1\,$\mu$g/m$^3$ for hypertension}, a monotone quartile gradient ($p_{\text{trend}}$ =
0.003) with the top quartile at 50\% greater odds, and \textbf{PAF 18.6\%}. \textbf{No
association with MetS overall or with any other component} --- the null is as informative
as the positive, and the authors report it rather than burying it. Three limitations bound
the inference: the design is cross-sectional so temporality is not established; exposure is
modelled at satellite resolution over a region with almost no reference monitoring and no
local validation is reported; and household air pollution from solid-fuel cooking, which
dominates personal exposure in this setting, is not in the model at all.}
{10.1016/j.envres.2026.125484}

\paperentry{Ting 2026}
{Statistical fragility and methodological constraints in metabolomic mediation of air pollution and cardiovascular risk}
{JACC Adv}
{\textbf{Metadata-only: correspondence deposited without an abstract.} Registered because it
completes an exchange this watch is already carrying --- the authors' reply
(\texttt{10.1016/j.jacadv.2026.103149}) surfaced in the 12 August window as an
Elsevier metadata-only record and remains unresolved. The subject, statistical fragility of
mediation estimates in air-pollution metabolomics, is directly relevant: mediation analysis
on high-dimensional metabolite panels with modest cohort sizes produces indirect effects
whose confidence intervals are dominated by the product-of-coefficients variance, and the
literature routinely reports them as though they were primary. No numeric claim is made from
this record. Retrieve both letters together once full text is available.}
{10.1016/j.jacadv.2026.103150}

\band{Burden, policy \& mitigation\hfill 3 records}{Coral}

\paperentry{Liu M.\ et al.\ 2026}
{Divergent near-surface ozone and NO$_2$ trends and ozone-attributable burden, Huaihe River Basin}
{Environ Monit Assess}
{Twelve years (2013--2024) of high-resolution satellite-derived surface MDA8 O$_3$ and
NO$_2$ across the Huaihe Basin, with pixel-wise OLS trends, the Hurst exponent as a
persistence diagnostic, and BenMAP-CE for attributable premature mortality and economic
loss. \textbf{O$_3$ rose at 2.792\,$\mu$g\,m$^{-3}$\,yr$^{-1}$ (mean
104.15\,$\mu$g\,m$^{-3}$) while NO$_2$ fell at $-$1.385\,$\mu$g\,m$^{-3}$\,yr$^{-1}$ (mean
29.86)} --- the textbook VOC-limited response to NOx-only control, here quantified at basin
scale with an outward expansion of O$_3$-affected area and contraction of NO$_2$ hotspots.
Cardiovascular mortality was the largest cause-specific component, concentrated in Jiangsu
and Shandong. \textbf{No PM metric is estimated}, which for a basin whose PM$_{2.5}$ and
O$_3$ share precursors is a real gap; and the Hurst exponent applied to a 12-point annual
series is a weak persistence estimator being asked to carry a forward-looking claim.}
{10.1007/s10661-026-15802-3}

\paperentry{Mou et al.\ 2026}
{Ship emissions and regional air quality in the post-DECA era, Yangtze River Basin}
{Environ Pollut}
{A bottom-up 2023 emission inventory for \textbf{18{,}478 vessels} built from
\textbf{$\approx$230 million AIS trajectories} on the world's busiest inland shipping
corridor, coupled to WRF-Chem for concentration attribution. SO$_2$ and PM$_{2.5}$ ship
emissions are substantially controlled post-DECA, but \textbf{NOx reached
156.41\,kt\,yr$^{-1}$, 82.0\% of major-pollutant mass}, with hotspots at
Yichang--Jingzhou, Wuhan and Nanjing--Shanghai. Ship contribution to NO$_2$ was
\textbf{0.16--3.41\,$\mu$g\,m$^{-3}$} annual mean across 11 mainstem cities; PM$_{2.5}$
increments were small but nitrate-dominated; O$_3$ fell in January and October under
near-source NO titration and rose in July. The policy read is clean --- fuel-sulfur control
has done its work and the residual problem is NOx --- but the inventory rests on activity-based
emission factors with no plume measurement to validate them, and the PM$_{2.5}$ increment
excludes secondary nitrate formed downwind of the model domain.}
{10.1016/j.envpol.2026.128946}

\paperentry{Liu D.\ et al.\ 2026}
{COPD burden 1990--2050 with risk-factor attribution across age and SDI}
{Medicine (Baltimore)}
{GBD 2021 across 204 countries, analysed by joinpoint regression, age--period--cohort
models and decomposition into population growth, ageing and epidemiological change, with
Bayesian APC projection to 2050. Age-standardised rates fell (AAPC $-$0.05\% incidence,
$-$1.49\% mortality, $-$1.48\% DALYs) while absolute burden rose to \textbf{213.4 million
cases and 1.69 million deaths in 2021}. \textbf{Ambient PM$_{2.5}$ carries 15--25\% of
DALYs and household air pollution 0--55\%}, against smoking at 38--55\%, with an inverted-V
SDI relationship peaking at SDI 0.4--0.6. The household-air-pollution range spanning zero
to fifty-five percent is the whole story and is presented as a single interval; it is a
statement about between-region heterogeneity, not uncertainty, and the paper does not
separate the two. GBD risk-factor PAFs inherit their exposure--response functions from the
integrated exposure-response curves, so this is a re-analysis of GBD outputs rather than
independent evidence. Tier C accordingly.}
{10.1097/MD.0000000000050162}

\band{Reproductive \& developmental\hfill 2 records}{Sky}

\paperentry{Hung et al.\ 2026}
{Late-pregnancy air pollution, gestational hypertension and neonatal anthropometry}
{iScience}
{Women recruited in the second trimester and classified gestational-hypertensive or
normotensive, with PM$_{2.5}$, PM$_{10}$, NO$_2$, SO$_2$, CO and O$_3$ over the 30 days
before delivery assigned from ambient monitors, and maternal peridelivery hs-CRP, IL-1$\beta$
and TNF-$\alpha$ plus cord TNF-$\alpha$ as candidate mediators. GH women had higher hs-CRP
and delivered lighter neonates with smaller chest and abdominal circumferences; PM,
NO$_2$, SO$_2$ and CO were \textbf{inversely associated with birth length}. \textbf{Mediation
analysis found no significant indirect effect through any measured cytokine}, which is the
useful negative --- the inflammatory-mediation hypothesis is the default assumption in this
literature and it did not hold here. Ambient monitor assignment over a 30-day window is
coarse for a length endpoint measured in centimetres, no sample size or effect sizes appear
in the abstract, and a single peridelivery cytokine draw is a poor proxy for a
gestation-long inflammatory exposure.}
{10.1016/j.isci.2026.117029}

\paperentry{Wang et al.\ 2026}
{PM$_{2.5}$ modifies the temperature--semen quality association}
{Environ Pollut}
{Prospective cohort with repeated semen analysis, linear mixed models with natural cubic
splines for temperature, then piecewise regression stratified into higher and lower
PM$_{2.5}$ groups across the low- and high-temperature segments. Temperature showed
\textbf{inverted-U associations with semen quality over the 0--90, 0--9 and 10--14\,d
windows} --- 0--90\,d spans a full spermatogenic cycle, 10--14\,d the epididymal transit,
so the window structure is biologically coherent. \textbf{Synergy between high temperature
and PM$_{2.5}$, antagonism between low temperature and PM$_{2.5}$}, both acting mainly on
progressive and total motility. Stratifying at a PM$_{2.5}$ median is a crude interaction
test that discards within-stratum contrast, and antagonism at low temperature is at least as
easily explained by seasonal confounding (heating-season composition, behaviour, indoor
time) as by a joint biological mechanism. No effect sizes or intervals are reported in the
abstract, so nothing from this record enters the forest plot.}
{10.1016/j.envpol.2026.128910}

\band{Respiratory \& allergic\hfill 1 record}{Sage}

\paperentry{Shi et al.\ 2026}
{Air pollution and AECOPD hospitalisation across three Hebei cities}
{Medicine (Baltimore)}
{Ecological time series on \textbf{14{,}786 AECOPD hospitalisations} from 49 hospitals in
Shijiazhuang, Handan and Chengde, 2018--2020 (1{,}096 days), with DLNMs adjusted for
long-term trend, temperature, humidity, day of week, holidays and pandemic stage, then
pooled by random-effects multivariate meta-analysis. Cumulative effects strengthened with
lag length and peaked at lag 0--5\,d for every pollutant while single-day effects were weak.
\textbf{PM$_{2.5}$ at the 99th percentile versus median gave RR 1.391 (1.169--1.658) at lag
0--5}; all five non-ozone pollutants showed non-linear cumulative exposure--response, and
higher 8-h ozone was associated with \emph{lower} admission risk, which the authors do not
explain and which is most plausibly seasonal confounding. The abstract's own conclusions are
internally inconsistent on susceptibility --- the results state associations were more
pronounced in males, the conclusion says females --- so the stratified findings should not be
cited without the full text. Six correlated pollutants tested across six lag structures with
no multiplicity adjustment.}
{10.1097/MD.0000000000050080}

\band{Occupational \& indoor exposure\hfill 1 record}{Clay}

\paperentry{Ranjbar et al.\ 2026}
{1-Nitropyrene as a diesel-specific indoor exposure proxy in California homes}
{Preprints.org (preprint)}
{1-Nitropyrene and the isomeric 2-nitropyrene and 2-nitrofluoranthene measured in indoor air
and house dust from \textbf{40 homes} in the eastern San Francisco Bay Area across two
sampling events, alongside continuous indoor black carbon. The motivation is specificity:
BC, NOx and PM$_{2.5}$ are all used as diesel surrogates and none is diesel-specific.
Detection was \textbf{74\% in air and 97\% in dust}, medians \textbf{0.42\,pg/m$^3$ (IQR
0.32--0.54)} and \textbf{340\,pg/g (IQR 190--620)}; air--dust correlation $r$ = 0.40
($p$ = 0.02) and BC--1-NP $r$ = 0.38 (air) and 0.47 (dust). Diagnostic ratios
(2-NFl:1-NP $<$ 5) indicate primary emission rather than secondary formation. Dust 1-NP
associated with more spatial diesel predictors than air 1-NP did, consistent with dust
integrating longer-term deposition --- which is the actual argument for using it. Preprint,
not peer reviewed. Forty homes, two events, no personal sampling, and a picogram-per-cubic-metre
analyte whose blank and recovery performance is not summarised in the abstract.}
{10.20944/preprints202608.0873.v1}

\band{Other clinical endpoints\hfill 1 record}{Slate}

\paperentry{Mandal et al.\ 2026}
{Long-term ambient PM$_{2.5}$ and kidney function in urban Indian adults}
{Kidney Int Rep}
{Population-representative cohort of \textbf{12{,}271 adults} aged 20 and over in Chennai and
Delhi, with eGFR (CKD-EPI 2009) at three points between 2010 and 2016 and annual mean
PM$_{2.5}$ assigned at residential geocode from a 1\,km national daily spatiotemporal model.
City-specific linear mixed models with inverse-probability weighting for loss to follow-up,
plus effect modification by sex, waist--hip ratio, hypertension and diabetes. Exposure
contrast between cities is extreme --- Chennai 33/37/30 versus Delhi 118/123/130\,$\mu$g/m$^3$
--- and both cities show decline: \textbf{$-$0.32 (95\% CI $-$0.49 to $-$0.15)} in Chennai and
\textbf{$-$0.42 ($-$0.57 to $-$0.27)\,ml/min/1.73\,m$^2$ per 5\,$\mu$g/m$^3$} in Delhi. That
the two cities agree despite a fourfold concentration difference is reassuring for the shape
of the relationship and awkward for its magnitude. CKD-EPI 2009 is poorly calibrated in South
Asian populations, three measurements over six years is thin for an eGFR slope, and the
by-city stratification means the between-city contrast --- the largest exposure contrast in the
data --- is never used.}
{10.1016/j.ekir.2026.106693}

\clearpage

% ==================================================================== TABLE
\section*{\sffamily\bfseries\color{Deep}Corpus register}
\vspace{-1.5mm}{\color{Amber}\rule{\linewidth}{1.1pt}}\vspace{2.5mm}

{\fontsize{6.7}{8.4}\selectfont
\setlength{\tabcolsep}{2.4pt}
\begin{longtable}{@{}>{\RaggedRight}p{25mm}>{\RaggedRight}p{41mm}>{\RaggedRight}p{27mm}%
>{\RaggedRight}p{23mm}>{\RaggedRight}p{18mm}>{\RaggedRight}p{15mm}@{}}
\toprule
\textsf{\textbf{First author}} & \textsf{\textbf{Focus}} & \textsf{\textbf{Journal}} &
\textsf{\textbf{Design}} & \textsf{\textbf{Metric}} & \textsf{\textbf{Region}}\\
\midrule
\endfirsthead
\toprule
\textsf{\textbf{First author}} & \textsf{\textbf{Focus}} & \textsf{\textbf{Journal}} &
\textsf{\textbf{Design}} & \textsf{\textbf{Metric}} & \textsf{\textbf{Region}}\\
\midrule
\endhead
\multicolumn{6}{@{}l@{}}{\cellcolor{Mist}\textsf{\textbf{\textcolor{Deep}{Neuro / mental health}}}}\\[0.6mm]
\href{https://doi.org/10.1001/jamanetworkopen.2026.25015}{Dales} & Wildfire-sourced PM$_{2.5}$ and ED visits for migraine and other primary headache & \textit{JAMA Netw Open} & Case-crossover & PM$_{2.5}$ (wildfire) & Canada\\
\href{https://doi.org/10.1093/ageing/afag215}{Glover} & Umbrella meta-review of urban environmental exposures and cognitive health & \textit{Age Ageing} & Umbrella review / meta-analysis & PM$_{2.5}$ & Global\\
\href{https://doi.org/10.1016/j.nsa.2026.107019}{Goossens} & In utero air pollution exposure and pre-/postnatal brain development & \textit{Neurosci Appl} & Narrative review & PM$_{2.5}$ / PM$_{10}$ / PM$_{\mathrm{coarse}}$ & Global\\
\href{https://doi.org/10.1016/j.isci.2026.116617}{Ma} & PM$_{2.5}$-induced neurotoxicity and depression: exposure risks and mechanisms & \textit{iScience} & Narrative review & PM$_{2.5}$ + components & Global\\
\href{https://doi.org/10.3760/cma.j.cn112150-20250804-00752}{Pan} & Maternal PM exposure and offspring neurodevelopment: effects and mechanisms & \textit{Zhonghua Yu Fang Yi Xue Za Zhi} & Narrative review & PM (bulk) & China\\
\addlinespace[0.9mm]
\multicolumn{6}{@{}l@{}}{\cellcolor{Mist}\textsf{\textbf{\textcolor{Teal}{Cardiovascular \& metabolic}}}}\\[0.6mm]
\href{https://doi.org/10.3760/cma.j.cn112338-20251120-00843}{Huang} & PM$_{2.5}$ and dyslipidaemia in rural Chinese children; dietary effect modification & \textit{Zhonghua Liu Xing Bing Xue Za Zhi} & Cross-sectional (multistage) & PM$_{2.5}$ & China\\
\href{https://doi.org/10.32394/pe/218966}{Kazmierska} & Air pollution and cardiovascular disease: review of current data & \textit{Przegl Epidemiol} & Narrative review & PM$_{2.5}$ / PM$_{10}$ & Global\\
\href{https://doi.org/10.1007/s00125-026-06814-2}{Xi} & Air pollution, built environment and temperature vs. prediabetes phenotypes & \textit{Diabetologia} & Prospective cohort & PM$_{2.5}$ & Germany\\
\href{https://doi.org/10.1161/JAHA.125.046917}{Yang} & Short-term PM$_{2.5}$ and hospitalisation/ED use in atrial fibrillation & \textit{J Am Heart Assoc} & Case-crossover & PM$_{2.5}$ (wildfire-influenced) & USA\\
\addlinespace[0.9mm]
\multicolumn{6}{@{}l@{}}{\cellcolor{Mist}\textsf{\textbf{\textcolor{Sky}{Reproductive \& developmental}}}}\\[0.6mm]
\href{https://doi.org/10.1016/j.jes.2025.10.043}{Chen} & Prenatal air pollution, fetal growth trajectory phenotypes and child neurodevelopment & \textit{J Environ Sci (China)} & Prospective cohort & PM$_{2.5}$ / PM$_{10}$ & China\\
\href{https://doi.org/10.1016/j.ecoenv.2026.120566}{Huang} & Joint prenatal bisphenol + air-pollutant mixture exposure and preschool IQ & \textit{Ecotoxicol Environ Saf} & Prospective cohort (mixtures) & PM$_{2.5}$ / PM$_{10}$ & China\\
\href{https://doi.org/10.1007/s40572-026-00552-8}{Peterson} & Ambient air pollution, PAHs and pubertal development: critical review & \textit{Curr Environ Health Rep} & Critical review & PM$_{2.5}$ components & Global\\
\href{https://doi.org/10.1016/j.envres.2026.125313}{Xie} & Long-term PM$_{2.5}$ and timing of pubertal development in boys and girls & \textit{Environ Res} & Prospective cohort & PM$_{2.5}$ & China\\
\addlinespace[0.9mm]
\multicolumn{6}{@{}l@{}}{\cellcolor{Mist}\textsf{\textbf{\textcolor{Sage}{Respiratory \& allergic}}}}\\[0.6mm]
\href{https://doi.org/10.1007/s00484-026-03289-z}{Bayar Muluk} & Climate change and otolaryngology: airway disease impacts and emission reduction & \textit{Int J Biometeorol} & Narrative review & PM$_{2.5}$ & Global\\
\href{https://doi.org/10.4103/mgr.MEDGASRES-D-25-00164}{Liao} & Short-term air pollution, respiratory symptoms and lung function in asthma & \textit{Med Gas Res} & Panel study & PM$_{2.5}$ / PM$_{10}$ & China\\
\href{https://doi.org/10.3390/epidemiologia7040097}{Sanchez de Prada} & Ambient PM and asthma-coded admissions in children 0-3 y, Spain & \textit{Epidemiologia} & Ecological panel & PM$_{2.5}$ / PM$_{10}$ & Spain\\
\href{https://doi.org/10.1016/j.jhazmat.2026.143076}{Yang} & Allergenic vs. toxic risk of airborne seed hairs compared with ambient PM$_{2.5}$ & \textit{J Hazard Mater} & Field measurement & PM$_{2.5}$ + bioaerosol & China\\
\addlinespace[0.9mm]
\multicolumn{6}{@{}l@{}}{\cellcolor{Mist}\textsf{\textbf{\textcolor{Violet}{Mechanistic toxicology}}}}\\[0.6mm]
\href{https://doi.org/10.1016/j.jes.2025.12.074}{Ding} & PM$_{2.5}$ promotes vascular calcification via miR-27b-5p/HES1 signalling & \textit{J Environ Sci (China)} & In vitro + organoid + murine & PM$_{2.5}$ & China\\
\href{https://doi.org/10.1007/s00360-026-01698-5}{Maraschi} & Osmoregulatory disruption and metal bioaccumulation from settleable atmospheric PM & \textit{J Comp Physiol B} & Controlled exposure (ecotox) & Settleable APM & Brazil\\
\href{https://doi.org/10.1016/j.jes.2025.11.018}{Wu} & Macrophage NLRP3 inflammasome in real-ambient PM-induced hepatic glucose dysregulation & \textit{J Environ Sci (China)} & Real-ambient murine + co-culture & Real-ambient PM & China\\
\addlinespace[0.9mm]
\multicolumn{6}{@{}l@{}}{\cellcolor{Mist}\textsf{\textbf{\textcolor{Clay}{Occupational \& indoor}}}}\\[0.6mm]
\href{https://doi.org/10.3389/fmedt.2026.1796828}{Chen} & Nested suction instrument to cut intraoperative surgical-smoke PM$_{2.5}$ & \textit{Front Med Technol} & Bench / device experiment & PM$_{2.5}$ (surgical smoke) & China\\
\href{https://doi.org/10.1093/annweh/wxag055}{Cole} & Determinants of respirable crystalline silica exposure in Australian tunnelling & \textit{Ann Work Expo Health} & Observational exposure study & Respirable dust / RCS & Australia\\
\href{https://doi.org/10.1016/j.jhazmat.2026.143081}{Ding} & Stage-resolved hair Zn-isotope dynamics as a firefighting combustion-exposure tracer & \textit{J Hazard Mater} & Repeated-measures biomarker & Combustion-derived PM & China\\
\href{https://doi.org/10.3390/jox16040119}{Ramsden} & Xenobiotic hazards in aircraft cabin air & \textit{J Xenobiot} & Review + risk economics & Ultrafine particles & Global\\
\addlinespace[0.9mm]
\multicolumn{6}{@{}l@{}}{\cellcolor{Mist}\textsf{\textbf{\textcolor{Amber}{Sensing, forecasting \& instrumentation}}}}\\[0.6mm]
\href{https://doi.org/10.1016/j.jenvman.2026.130515}{Shi} & Adaptive spatiotemporal graph transformer for multi-site multi-step PM$_{2.5}$ forecasting & \textit{J Environ Manage} & Deep-learning model & PM$_{2.5}$ & China / India\\
\href{https://doi.org/10.3390/bios16070380}{Tsuchiya} & Seasonal and indoor-outdoor characteristics of fluorescent aerosol particles in offices & \textit{Biosensors} & Real-time sensor campaign & Bioaerosol / coarse & Japan\\
\addlinespace[0.9mm]
\multicolumn{6}{@{}l@{}}{\cellcolor{Mist}\textsf{\textbf{\textcolor{Amber}{Exposure assessment \& modelling}}}}\\[0.6mm]
\href{https://doi.org/10.1029/2025GH001621}{Solorzano-Araque} & Multiscale exposure modelling: LOTOS-EUROS + mobility surveys + particle toxicity & \textit{GeoHealth} & CTM + mobility integration & PM$_{2.5}$ / PM$_{10}$ & Colombia\\
\addlinespace[0.9mm]
\multicolumn{6}{@{}l@{}}{\cellcolor{Mist}\textsf{\textbf{\textcolor{Coral}{Burden, policy \& mitigation}}}}\\[0.6mm]
\href{https://doi.org/10.1002/hsr2.72886}{Karamzad} & Ischaemic heart disease burden in MENA 1990-2021; PM attributable fraction 28.2\% & \textit{Health Sci Rep} & GBD secondary analysis & Ambient PM & MENA\\
\href{https://doi.org/10.1371/journal.pone.0347989}{Rees} & HEPA filtration to reduce respiratory infection in care homes & \textit{PLoS One} & RCT process evaluation & Filtered PM & UK\\
\href{https://doi.org/10.1016/j.jes.2026.01.009}{Shan} & Policy-driven emission reductions and post-pandemic rebound, high-resolution inventories & \textit{J Environ Sci (China)} & Emission inventory & PM emissions & China\\
\addlinespace[0.9mm]
\multicolumn{6}{@{}l@{}}{\cellcolor{Mist}\textsf{\textbf{\textcolor{Slate}{Other clinical endpoints}}}}\\[0.6mm]
\href{https://doi.org/10.3390/dj14070440}{Lupita} & PM$_{2.5}$, preventive-care capacity and caries/edentulism burden across the EU & \textit{Dent J} & Ecological panel & PM$_{2.5}$ & EU-27\\
\href{https://doi.org/10.1016/j.jdin.2026.06.008}{Park} & Long-term PM exposure and risk of androgenetic alopecia & \textit{JAAD Int} & Nationwide cohort & PM$_{2.5}$ / PM$_{10}$ & South Korea\\
\href{https://doi.org/10.3171/2026.3.PEDS2615}{Sun} & Environmental exposure and area deprivation vs. paediatric hydrocephalus outcomes & \textit{J Neurosurg Pediatr} & Registry cohort + geospatial & PM$_{2.5}$ (satellite-derived) & USA\\
\addlinespace[0.9mm]
\bottomrule
\end{longtable}}


\vspace{4mm}
\begin{keybox}
{\sffamily\bfseries\footnotesize\color{Deep}Provenance and method}\\[1.2mm]
{\fontsize{7.2}{8.8}\selectfont
\textbf{Window.} PubMed \texttt{[EDAT]} \textbf{2026-08-15 to 2026-08-15}, contiguous with
the 14 August issue; no gap. Every record was checked against \texttt{state/seen.json}
(pmid $\rightarrow$ lowercased doi $\rightarrow$ title signature) before summarising;
\textbf{19 Europe PMC hits were dropped as already-seen duplicates} of records carried on
8--14 August, the largest cross-source collision so far.\\[0.8mm]
\textbf{Per-leg yield.} PubMed connector, two separate queries on the health and
instrumentation axes: the health axis returned 7 hits and the instrumentation axis, once
narrowed to exclude the \texttt{calibration}/\texttt{exposure model} MeSH explosions that
returned 76 mostly-clinical records, returned \textbf{1} --- fully contained in the health
set. \textbf{6 carried} from the connector. Local harvester: \texttt{pubmed\_health} 5, all
contained in the connector set; \texttt{pubmed\_sensing} 0. \textbf{Europe PMC carried the
issue}: 59 hits on \texttt{CREATION\_DATE}, 40 new after deduplication, \textbf{15 carried}
--- including every instrumentation record in this digest. Crossref by ISSN returned 9
across 18 journals, of which one (\textit{Air Qual Atmos Health}, North African urban
anomalies) was in scope but arrived after screening closed and is deferred to the weekly
sweep; the rest were \textit{Environ Sci Technol} and \textit{Environ Monit Assess} records
outside scope or already carried. arXiv 146\,kB, nothing on or after 11 August.
\textbf{OpenAlex HTTP 429, 13th consecutive} --- the leg is dead and no further effort is
warranted. ClinicalTrials.gov: no PM or air-pollution registration created or updated inside
the window; \texttt{trials.json} unchanged.\\[0.8mm]
\textbf{Harvester defect found.} \texttt{harvest.py}'s Europe PMC leg returned \textbf{0}
for this window while a hand-issued \texttt{CREATION\_DATE:[2026-08-15 TO 2026-08-15]}
query against the same endpoint returned \textbf{59}. Fifteen of today's 21 records came
from the manual query, so the committed leg would have shipped a six-record issue. Logged
to \texttt{PIPELINE.md} for repair before the next run; the daily harvest cannot be trusted
on the Europe PMC leg until it is fixed.\\[0.8mm]
\textbf{Rejections (26).} By reason: eleven keyword collisions where \texttt{aerosol} or
\texttt{air pollution} matched an unrelated field (PIPAC oncology, therapeutic aerosol in
equine wound care, free-space optical communication, solar UV dosimetry, veterinary
virology, CF pharmacology, infection-control adherence, broiler LCA, Nile Red photophysics,
X-ray thermometry of glycerol droplets, a thermoelectric gas sensor with no particulate
channel); four pharmaceutical-inhaler records; three emissions/chemistry papers with no PM
measurement (NOx budget, brominated dioxins, green nanomaterials); a VOC-only vehicle
catalyst study framed as PM-precursor work but reporting no PM; a scoping review whose
literature search closed in April 2021; an allergic-rhinitis DLNM carrying meteorology only,
with pollution as an adjustment covariate; a stakeholder-perception window-design study with
no measured indoor PM; a veterinary climate-health commentary; a duplicate methodological
correspondence; and one PMC front-matter item. All logged with reasons to
\texttt{state/rejected.jsonl}.\\[0.8mm]
\textbf{Two records are metadata-only} (Chen et al., \textit{Allergy}; Ting,
\textit{JACC Adv}) --- deposited without abstracts, registered by title only, and supporting
no numeric claim anywhere in this issue. \textbf{One record is a preprint} (Ranjbar et al.,
Preprints.org) and is labelled as such in the register.\\[0.8mm]
\textbf{Labels.} Subtopic, design, particle-metric and geography labels were drawn only from
vocabularies already mapped in \texttt{plots.py} and \texttt{mktable.py}; no new leaf label
was written this issue. Assignment is single-label by dominant contribution and is a
judgement, not a controlled vocabulary. \texttt{check\_dois.py} returned 0 fail, 0 warn
across all 21 records.\\[0.8mm]
Content derived from PubMed, Europe PMC and Crossref metadata; abstracts remain the
copyright of their respective publishers. All DOI links resolve to the publisher of record.\par}
\end{keybox}

\end{document}
